\section{Subarray Sums I}
\label{sec:cses-subarray-sums-i}

\problemstatement{Given an array of $n$ \emph{positive} integers and a
target $x$, count the number of contiguous subarrays whose sum equals $x$.}

\begin{patternbox}
With all values positive, the window sum is \textbf{monotonic}: extending
the right end only increases it, and advancing the left end only decreases
it. That monotonicity is exactly what lets a \textbf{two-pointer} sliding
window count matches in one linear pass.
\end{patternbox}

\subsection*{Why two pointers work here}

Keep a window $[l,r]$ and its sum. Advance $r$, adding $a[r]$. Whenever the
sum exceeds $x$, advance $l$, subtracting $a[l]$ -- safe because positivity
guarantees shrinking only lowers the sum. Whenever the sum equals $x$,
record one subarray. Because the sum never decreases when $r$ moves right
and never increases when $l$ moves right, each pointer only moves forward,
giving $O(n)$.

\begin{ideabox}[Positivity Makes the Window Monotone]
Two pointers require that growing the window changes the tracked quantity
in one direction only. Positive values guarantee that for sums -- which is
why Subarray Sums~II (negatives allowed) must abandon two pointers for
prefix sums.
\end{ideabox}

\subsection*{Algorithm}

\begin{enumerate}
  \item $l=0$, $\textit{sum}=0$, $\textit{count}=0$.
  \item For $r=0\ldots n-1$: add $a[r]$; while $\textit{sum}>x$ advance
        $l$ subtracting $a[l]$; if $\textit{sum}=x$ increment
        $\textit{count}$.
  \item Output $\textit{count}$.
\end{enumerate}

\subsection*{Dry run}

\dryrun{$a=[2,4,1,2,7]$, $x=7$.}

\begin{itemize}
  \item $[2,4,1]$ sums to $7\Rightarrow$ count $1$.
  \item Window shrinks/extends; $[4,1,2]=7\Rightarrow$ count $2$; later
        $[7]=7\Rightarrow$ count $3$.
  \item Answer $3$.
\end{itemize}

\subsection*{C++ implementation}

\begin{lstlisting}
#include <bits/stdc++.h>
using namespace std;

int main() {
    int n; long long x;
    cin >> n >> x;
    vector<long long> a(n);
    for (auto& v : a) cin >> v;

    long long sum = 0, count = 0;
    int l = 0;
    for (int r = 0; r < n; r++) {
        sum += a[r];
        while (sum > x) sum -= a[l++];         // positivity: safe to shrink
        if (sum == x) count++;
    }
    cout << count << "\n";
    return 0;
}
\end{lstlisting}

\begin{complexitybox}
\textbf{Time Complexity.} $O(n)$ -- both pointers only move forward.
\textbf{Space Complexity.} $O(1)$.
\end{complexitybox}

\begin{takeawaybox}
Two pointers count target-sum subarrays in linear time precisely because
positive values keep the window sum monotone. The moment negatives enter,
monotonicity breaks and the prefix-sum-plus-hash-map method takes over.
\end{takeawaybox}
